\section{Parameters, ranks, and probability scales}
\label{sec:parameters}

This Appendix checks that the construction parameters can be chosen
in the required order and that the resulting bounds suffice at every
probability scale. All construction constants are fixed independently
of \(n\) and of the unit law. These constants may be very large.
The mixer matrices are then chosen separately at each sufficiently
large \(n\), by \Cref{lem:generic-mixers}, before the unit law is selected.

\subsection{The acyclic order of construction}

The construction constants and first pin budget from
\eqref{eq:early-constants} and \eqref{eq:pin-budgets} are
\begin{align}
 C_0&=1000,& g&=10^9+1,&
 D&=4C_0g,& M&=2^{1000},\nonumber\\
 k_{\max}&=56(g-1),&
 \zeta&=\frac{1}{1000k_{\max}},&
 K_1&=\left\lceil\frac{4(D+1)}{\zeta}\right\rceil .
 \label{eq:ledger-early}
\end{align}
In particular \(g\) is odd and \(D=4000g\).
The early law and repinning budgets are
\begin{align}
 r&=2K_1+2,&
 L_0&=\lceil100(D+10)\rceil,&
 d_0&=4rL_0,\nonumber\\
 u_0&=10(K_1+d_0+1),&
 K&\in\N,\qquad
 K\ge\frac{10(D+K_1+u_0+10)}{\zeta}.
 \label{eq:ledger-law}
\end{align}
These are the bounds in \eqref{eq:law-budgets}; the construction
takes \(K\) to be the least integer meeting the displayed lower bound.

The remaining early tolerances in
\Cref{thm:phase-alternative,lem:injection} include
\begin{equation}
 c_s=10^{-12}(1+r)^{-4},\qquad
 q_0=2^{-100r-10},\qquad
 \delta_{\mathrm{cov}}=.0005.
 \label{eq:ledger-cover-tolerances}
\end{equation}
For each \(n\), the moment order is the largest power of two
\(s\le c_sN\); it is even for sufficiently large \(n\), and
\(s\ge c_sN/2\).
Choose \(p_0>0\) sufficiently small in terms of \(r,D,c_s,q_0\)
and the fixed phase tolerance. Fix the transpose-injection error
thresholds as well, with their total over all component and endpoint
roles below the required constant error, such as \(.001\).
These thresholds precede selector and channel sizes.

The exponential demands of the no-cover argument are simultaneously
satisfiable in this order. Fix a phase-test threshold \(\vartheta>0\)
smaller than the required fixed phase accuracy, and choose a constant
\(\Gamma>D+2r+1\) large enough to absorb the indicated density and
bounded-rank tensor counts. With \(L_\vartheta=\log_2(1/\vartheta)\),
it suffices to arrange
\begin{align}
 \frac{c_s}{2}
 \left(\frac{q_0}{2}\log_2(1/p_0)-2-L_\vartheta\right)
     &>\Gamma,\label{eq:ledger-p0-demand}\\
 \frac{c_s}{2}
 \left(\frac{hq_0}{4}-L_\vartheta\right)
     &>\Gamma.\label{eq:ledger-h-moment-demand}
\end{align}
The first determines a sufficiently small \(p_0\) using only early
quantities; the second is a later lower bound on \(h\).
Indeed the two moment terms are
\(2^{2s}p_0^{q_0s/2}\) and \(2^{-hq_0s/4}\), and Markov at threshold
\(\vartheta\) costs \(2^{sL_\vartheta}\).
The comparison between raw and independent channel laws has a factor
depending on \(h\), but it is constant in \(n\) and does not alter the
linear exponents.

Next choose the row and selector parameters:
\begin{align}
 R_*&=2(K+20),&
 j_*&\ge10(K+1)(R_*+20),\nonumber\\
 B_0&=2|\cE|K(R_*+14),&
 B_*&=10(B_0+10000g+1),\nonumber\\
 2^{b-2j_*}&>100g(B_*+1),&
 p_*&=\sum_{j=0}^{j_*}\binom bj,\nonumber\\
 A&=10(K+1)p_*(g^2+5),&
 s_0&=4\lceil A\rceil,\qquad
 J>100(s_0+K^2+1).
 \label{eq:ledger-selectors}
\end{align}
Integer inequalities are met by rounding upward when needed.
This agrees with \Cref{lem:label-support,lem:generic-mixers}.
In particular \(j_*\) exceeds the interpolation degree required for
\((K+1)(R_*+14)\) labels. The length \(b\) is chosen before \(A\),
and hence before \(s_0,J\).

Now choose a sufficiently large integer \(r_0\), always setting
\begin{equation}
                         h=1000r_0.
 \label{eq:ledger-channels}
\end{equation}
The demands are finite in number and depend only on preceding
constants:
\begin{enumerate}
\item the moment demand \eqref{eq:ledger-h-moment-demand};
\item fixed-cover and fresh-probe transpose-injection errors in
  \Cref{thm:phase-alternative};
\item the accepting-set family estimate in \Cref{lem:injection};
\item residual ranks and private-channel space in
  \Cref{thm:gradient-realization};
\item the tester rank inequality in \Cref{lem:zero-delta}.
\end{enumerate}
Each is eventually satisfied by increasing \(r_0\).
Quantitative checks appear below.

Finally choose a sufficiently large integer \(M_0\), and put
\begin{equation}
                  N=M_0n,\qquad m=2^{C_0gN}.
 \label{eq:ledger-final-dimensions}
\end{equation}
The coefficient of every \(O(n)\) count paid as a fraction of \(N\)
has then been fixed. Choose \(M_0\) to make all those fractions
smaller than their prescribed slacks. Only afterward does \(n\)
tend to infinity through integers with \(n\ge2r_0\).
The resulting dependency order is
\[
 \begin{gathered}
 (C_0,g,D,M,k_{\max},\zeta,K_1)
 \longrightarrow
 (r,L_0,d_0,u_0,K,\text{early tolerances},p_0)\\
 \longrightarrow
 (R_*,j_*,B_0,B_*,b,p_*,A,s_0,J)
 \longrightarrow(r_0,h)
 \longrightarrow M_0
 \longrightarrow n .
 \end{gathered}
\]

\subsection{The channel demands precede later tests}

In the frequent-output argument, a fixed collection of \(L_0\)
fresh probes succeeds on B-mass at least
\begin{equation}
       q_*=(\delta_{\mathrm{cov}}/2)
              (\delta_{\mathrm{cov}}/4)^{L_0}>0.
 \label{eq:ledger-probe-mass}
\end{equation}
The one-vertex density factor depends only on the early marked-law
restrictions and \(p_0\).
Transpose-injection failure on the fixed fresh vectors is at most
that factor times \(O(L_0)2^{L_0-h}\).
It can be made smaller than \(q_*/2\) by increasing \(h\), without
consulting a selector choice or a future common test.
The later raw-event exponent is at most
\[
       -(.99-.51)L_0N+DN+\text{small linear errors}.
\]
The specified \(L_0\) leaves a margin close to \((.48L_0-D)N\).
The small coefficient-count errors are paid by \(M_0\).

For \Cref{lem:injection}, one may use accepting-set log-size
\begin{equation}
                       C_sN,\qquad C_s=100(r+d_0+1).
 \label{eq:ledger-accepting-family}
\end{equation}
The family is encoded by rank-bounded \(\Delta\)'s and channel
functionals on a cover tensor space of dimension at most \(d_0N\).
There is no \(h\) multiplier in its \(N\)-coefficient.
Let \(c>0\) be the fixed per-probe success lower bound after the early
failure tolerances are fixed.
With \(\ell=\lfloor N/(10(K+1))\rfloor\), the comparison is
\[
       c^\ell
       \quad\hbox{against}\quad
       2^{-(h-2K)\ell+o(N)}2^{C_sN}.
\]
For example it suffices that
\begin{equation}
 h>2K+\log_2(1/c)+10(K+1)(C_s+2).
 \label{eq:ledger-injection-demand}
\end{equation}
The floor changes only a bounded factor. Every quantity on the right
was fixed before \(h\).

For channel realization, the explicit bounds in
\Cref{eq:baseline-rank-bound,eq:residual-rank-bound,eq:block-compression-bound,eq:quotient-compression-bound}
are
\[
\begin{aligned}
 B_{\mathrm{lin}}&=3K+28,&
 R_{\mathrm{res}}&=14K+140,\\
 D_{\mathrm{blk}}&=2p_*|\cE|(K+14+R_{\mathrm{res}}),&
 D_{\mathrm{quo}}&=2p_*\bigl(1+(g^2+3)D_{\mathrm{blk}}\bigr)+2K.
\end{aligned}
\]
The pure correction rank per component is at most
\(30r_0+28J|\cE|+D_{\mathrm{quo}}\), and the remaining channel
pairing has rank at least \(h-2K-4B_{\mathrm{lin}}\).
Thus \Cref{eq:realization-r0-requirement} asks precisely for
\begin{equation}
       970r_0>28J|\cE|+D_{\mathrm{quo}}+2K+4B_{\mathrm{lin}}.
 \label{eq:ledger-realization-demand}
\end{equation}
Every quantity on the right is independent of \(h,r_0,n\)
and is fixed by parameters through \(J\).
The baseline first factors through bounded pin/key projections on
primal inputs, then derivative corrections of bounded rank are
chosen, and the small residual is compressed before the
\(r_0\)-sized pure terms are factored through private channels.

The tester inequality in \Cref{lem:zero-delta} is
\begin{equation}
 \frac{g-1}{2}\left(g-\frac{4D}{g}\right)r_0
          >2(g-1)(h+2JK_1)+1.
 \label{eq:ledger-pure-tester-demand}
\end{equation}
Substituting \(D=4000g\) and \(h=1000r_0\) gives
\[
       \frac{g-20000}{2}\,r_0
             >4JK_1+\frac1{g-1}.
\]
The coefficient on the left is positive for the specified \(g\).
This is therefore another compatible lower bound on \(r_0\).

\subsection{Pin, atom, and tensor-rank ledger}

\begin{center}
\begin{tabular}{@{}p{.39\textwidth}p{.53\textwidth}@{}}
\toprule
Quantity & Bound and role\\
\midrule
First pin dimension & \(K_1\), summed over components and signs.\\
Added exact support pins & At most \(4K_1\), for both old marked
tensors in a unit.\\
Added cover-evaluation pins & At most \(2d_0\), for both endpoints'
channel products on cover bases.\\
Total added rank & \(4K_1+2d_0<u_0\).\\
Final pin dimension & At most \(K\); the second peeling is needed
only in case \textup{(iii)} of \Cref{prop:status-preparation}.\\
Effective space & \(\dim C_i\le K^2\); every member has total component
rank at most \(K\).\\
Selected atoms at one endpoint & At most \(14\), across its two lists.\\
Component rank in the annihilator argument &
\(2(K+14)\le R_*=2(K+20)\).\\
List length per cross edge & At most \(7\) on either side.\\
Atom positions in one unit's full query & At most \(4\cdot7=28\).\\
General key-vector slots &
\(28\cdot2(g-1)=k_{\max}=56(g-1)\).\\
Rare-query key-vector slots & \(4\cdot2(g-1)=8(g-1)\).\\
\bottomrule
\end{tabular}
\end{center}

Each tag atom occupies \(g-1\) star components and supplies one plus
and one minus vector at each. These are vector counts, not scalar
Gram counts.
The label exclusions in \Cref{lem:sparse-pin-labels} give nominal
independence modulo the pins, as proved in
\Cref{lem:small-tables}, part~\textup{(ii)}.
The effective-space dimension is counted in projected pin bases:
for a fixed endpoint \(\sum_e\dim S_{i,e}^\pm\le K\), so
\[
 \sum_e\dim(S_{i,e}^+\otimes S_{i,e}^-)
 \le
 \left(\sum_e\dim S_{i,e}^+\right)
 \left(\sum_e\dim S_{i,e}^-\right)
 \le K^2.
\]
Thus an effective tensor basis is encoded in bounded projected
coordinates, not as unrestricted matrices with \(O(n^2)\) entries.

The unprojected representation of one gradient in
\Cref{thm:gradient-realization} has per-component rank at most
\[
                         15r_0+14J|\cE|.
\]
Two endpoint contributions give \(30r_0+O(1)\).
If \(A_+,A_-\) remove bounded pin/key subspaces, then
\[
 \rank(M-A_+\trans M A_-)
       \le\rank(I-A_+)+\rank(I-A_-).
\]
Hence this projection contributes only bounded residual rank,
without multiplying the tester rank by a pin-dependent constant.
The later blockwise compression and derivative corrections preserve
that independence from \(r_0\).

\subsection{Scalar records and the final dimension multiplier}

Write
\begin{equation}
                  d_n=\dim\cB=p_*\bigl(1+(g^2+3)n\bigr).
 \label{eq:ledger-column-dimension}
\end{equation}
Then
\[
       \frac{d_n}{N}
       =\frac{p_*(g^2+3)}{M_0}+\frac{p_*}{M_0n}.
\]
This ratio, and the ratio after adding the bounded channel dimension,
can be made arbitrarily small by choosing \(M_0\) last.
In particular the image estimate can use a fixed sufficiently small
\(\varepsilon<\zeta/4\). The same choice leaves room for the injective
frames and prescribed self-Gram values.

The collision records are linear in \(n\).
There are at most \(56\) point inputs across two units, costing at most
\[
                         56(g^2+3)n+56b
\]
bits. A list of \(K\) nominal pin directions uses at most
\(2K(d_n+h)\) coefficient bits, plus bounded component, sign and rank
indices.
Once pin projections are fixed, a basis of \(C_i\) uses at most
\(K^4\) binary coefficients in its tensor-coordinate space;
there are four endpoints.
Tester summaries, numerical small tables, projection relationships
and finite statuses have bounded lengths in \(n\).
The construction matrices \(E,L,R\) are fixed data and are not
encoded afresh in each query record.

The total frozen dimension of a unit is at most \(K+k_{\max}\),
over all components and signs.
There are \(2(d_n+h)\) nominal columns of a given sign at each
component. Recording all interactions of both units' columns with
opposite frozen images therefore uses at most
\[
                         4(d_n+h)(K+k_{\max})
\]
scalar bits.
The deterministic solution is selected from these data and the point
parameters. No unrestricted \(O(n^2)\)-entry matrix is appended.
The number of record cells, including bounded nominal choices when
needed, is consequently at most
\begin{equation}
                         2^{C_{\mathrm{rec}}n+C_{\mathrm{rec},0}},
 \label{eq:ledger-record-count}
\end{equation}
for constants fixed before \(M_0\).

Only mixed primal-channel Gram entries are tested in the final target.
For one cross Gram block and component, the two endpoint primal
summands have total dimension \(2d_n\) and the channel summands have
total dimension \(2h\) on either side. The two mixed products
contain at most \(8d_nh\) entries.
There are two cross Gram blocks; hence at most
\begin{equation}
                              16|\cE|\,h\,d_n
 \label{eq:ledger-tested-gram-bits}
\end{equation}
bits are tested. Frozen or redundant entries only reduce this count.
This is \(O(n)\), with coefficient already fixed.

Choose \(M_0\) to meet every earlier coefficient/codimension slack
and, in particular, so that for sufficiently large \(n\),
\begin{equation}
 \log_2(\text{record-cell count})<.005N,\qquad
 \text{tested scalar Gram bits}<.01N.
 \label{eq:ledger-collision-bit-slacks}
\end{equation}
The smaller image and pin-count slacks form a finite collection
and can be met by the same final enlargement of \(M_0\).

The remaining collision estimates are those of
\Cref{prop:collision}. In particular, \eqref{eq:cell-pmax} gives
the following conditional point-mass bound on a surviving cell:
\[
 \begin{split}
 \log_2 p_{\max}/N
  &\le -(1-2\zeta)(k+t)+(k+.01)\\
  &=-(1-2\zeta)t+.01+2\zeta k
       \le-.95t
       \qquad(t\ge1,\ k\le k_{\max}).
 \end{split}
\]
Here \(2\zeta k\le.002\).
The quotient-character estimate \eqref{eq:gram-character-bound}
bounds the absolute mean of a rank-\(t\) character by
\[
       2^{tN/2}
       \bigl(2^{-.95tN}2^{-.95tN}\bigr)^{1/2}
             =2^{-.45tN}.
\]
This pays for fewer than \(.01N\) scalar bits, leaving agreement
probability at least \(2^{-.02N}\).
The overlap loss in \eqref{eq:fubini-pruning}, with the record count
made explicit, is at most
\[
 2|\cK_n|\,2^{C_{\mathrm{rec}}n+C_{\mathrm{rec},0}}
       2^{-(k+.01)N},
\]
which is exponentially small since \(|\cK_n|\le2^{kN}\) and
\eqref{eq:ledger-collision-bit-slacks} holds.
As in that proof, this bound is obtained at each fixed parameter
pair and then averaged.

\subsection{Probability scales and later analytic choices}
\label{sec:probability-scales}

The mixed law retains marginal control, whereas a normalized leaf
is required to satisfy a joint cap and the exact-image bounds.
Since \(M_0\) is fixed, \(o(N)=o(n)\) in all the comparisons below.

\begin{center}
\begin{tabular}{@{}p{.36\textwidth}p{.56\textwidth}@{}}
\toprule
Mechanism & Scale and permissible cost\\
\midrule
Original unit law & Joint cap \(2^{DN}\mu^2\); both marginals
at most \(M\mu\).\\
Prepared mixed law & Joint cap \(2^{(D+o(1))N}\mu^2\), marginals
at most \(M2^{o(N)}\mu\); restriction mass \(2^{-o(N)}\).\\
Prepared normalized leaves & Uniform joint cap
\(2^{C_{\mathrm{leaf}}N}\mu^2\) and the fresh-image bound with
\(2\zeta\) slack.\\
Fixed-cost early restrictions & Negative logarithmic mass cost
\(O(1)\), absorbed by fixed linear exponent margins.\\
Polynomial bad-law restriction & \(n^p=2^{o(n)}=2^{o(N)}\) for
fixed \(p\); mixed caps survive, but old separate leaf caps need not.\\
Single-endpoint typicality & Error \(2^{-\Omega(n)}\), surviving
\(2^{o(n)}\) marginal inflation.\\
Paired oversampling and cross-rank failures & Error
\(2^{-\Omega(n^2)}\), surviving \(2^{O(N)}=2^{O(n)}\) joint inflation.\\
Tiny-cover compatible pairs & Mass \(\Omega(N^{-4})\), with
exponential singleton-character and peeling errors.\\
Large-set hit exceptions & \(o(n^{-p})\) for every fixed \(p\),
uniform over deterministic large key sets and finite requests.\\
Independent B-leaf net use & Exception at most
\(2CL e_n/a_0=o(N^{-4})\), with \(C,L,a_0^{-1}\) fixed.\\
Key equality & Factor \(|\cK_n|^{-1}\ge2^{-kN}\),
\(k\le k_{\max}\).\\
Final four-hole probability &
\(2^{-(k_{\max}+.03)N-o(N)}\), exceeding \(2^{-100gN}\) for large \(n\).\\
\bottomrule
\end{tabular}
\end{center}

The polynomial compatibility scale is explicit.
The cover tensor space has dimension at most \(d_0N\), so
\(\rank(1+F)\le R_N:=1+d_0N\).
An incompatible clique has size at most \(R_N^2\).
Sampling \(R_N^2+1\) times and union bounding compatible pairs gives
probability at least \(\binom{R_N^2+1}{2}^{-1}\).
The four requested phase bits retain one quarter of this contribution,
up to exponentially small error.
A fixed fourth power therefore suffices; no exponent depending on
\(n\) is concealed in the term \emph{inverse-polynomial}.

Subexponential restrictions are compatible with leaf trimming.
If retained total mass is \(q_n=2^{-o(N)}\), leaves whose surviving
fraction is below \(2^{-\zeta N}\) contribute at most
\[
                     2^{-\zeta N}/q_n
                       =2^{-(\zeta-o(1))N}
\]
of the retained law, negligible compared with \(N^{-4}\).
The second-peeling threshold also loses exponentially little:
there are at most \(2^{u_0N}\) image choices and \(2^{O(n)}\)
other choices, while each discarded part has relative mass below
\(2^{-(u_0+1)N}\).

Later entropy bounds, nets, partitions, weak product accuracies and
finite transformation sample sizes are analytic choices made after
the construction constants. They may depend on \(M_0\), on a fixed
positive test-cell mass, or on fixed accuracy.
They do not change a vector-pin rank or the scalar record coefficient
in \eqref{eq:ledger-record-count}.
For each law sequence satisfying the stated density bounds and each
fixed finite test, these choices affect only constant factors and the
required onset in \(n\). In a countable common-test construction,
each fixed finite stage is passed before taking a diagonal subsequence.
No exponential rate uniform over all stages is asserted.

For example, unary positivity fixes
\(\gamma=2^{-K^2-3}\) and
\(b_{\mathrm{int}}=2^{-J+2s_0}<\gamma\).
On a reference key cell of mass at least \(\delta>0\), the proof
first obtains conditional input mass at least \(\delta/2\).
It may therefore use
\[
                    u>\frac{2}{\delta(\gamma-b_{\mathrm{int}})^2}
\]
transformations. This changes the constant in
\(O_u(2^{-s_0n})\), while \(s_0=4\lceil A\rceil\) and the profile
count \(2^{An}\) remain fixed. A smaller cell does not require a
larger \(J\).
Likewise the final B-net accuracy in \Cref{prop:rare-overlap} is fixed
after the positive density, set-size, flag-mass and A-hit thresholds.
Its cardinality is then a constant that a superpolynomial error absorbs.

The final collision margin is
\[
                    100g-[56(g-1)+.03]=44g+55.97>0.
\]
For sampling, the fingerprint log-count is
\[
                O(N^3 2^{100gN})=o(m),\qquad m=2^{1000gN}.
\]
For \(k=\lceil m/200\rceil\), the exceptional-pair bound is
\[
                         m^{2k}2^{-kDN}=2^{-2000gNk}.
\]
The terminal endpoint-set mass is below \(1/M=2^{-1000}\);
its binomial tail at \(m/200\) is exponentially small in \(m\).
These estimates pay for the fingerprint union in
\Cref{prop:raw-to-graph} and complete the scale ledger.
